\documentclass[11pt]{article}
\usepackage{amsmath,amssymb,amsthm,hyperref}
\newtheorem{lemma}{Lemma}
\newcommand{\E}{\mathbb E}
\title{An integrated polarization bound without exponential coefficient loss}
\author{Independent asymptotic companion calculation}
\date{30 September 2026}
\begin{document}\maketitle
Use the notation of \path{forced_success_mixture_comparison.tex}:
$\ell=\lfloor m/4\rfloor$, $N=m-\ell$, $\alpha=\ell/m$,
$\delta_j=1-p_j$, $s=\sum\delta_j$, $x_j=\ell\delta_j$,
$A_B=\prod_{j\in B}p_j$, and $x_R=N^{-1}\sum_{j\in R}x_j$ for
$R=B^c$. We use the established estimates
\[
 \max_j\delta_j\le256(s+1)/m,\qquad
 \E_B A_B\le e^{-\alpha s},\qquad
 \E_w e^{-cs}\le C_c/m \quad(c>0\text{ fixed}).        \tag{1}
\]
Fix $0<\theta\le1/16$. The polynomials $R_d$ and $P_d$ below are the
scaled Laguerre taper and normalized square from that companion note.
Their degrees are at most $2d-1$, and we require $2d\le N$.

\begin{lemma}[Centered elementary symmetric coefficients]
If $y_1,\ldots,y_N$ are real, $\sum y_j=0$, and $|y_j|\le D$, then
\[
 \frac{|e_k(y)|}{\binom Nk}
 \le\left(\sqrt e D\sqrt{k/N}\right)^k
 \qquad(1\le k\le N).
\tag{2}
\]
\end{lemma}
\begin{proof}
The elementary inequality
$|1+u|\le\exp(\operatorname{Re}u+|u|^2/2)$ gives, for any complex $z$,
\[
 \left|\prod_j(1+y_jz)\right|
 \le\exp\left(|z|^2\sum_jy_j^2/2\right).
\]
Cauchy's coefficient estimate at radius
$\sqrt{k/\sum y_j^2}$ yields
$|e_k(y)|\le(e\sum y_j^2/k)^{k/2}$; the case of zero sum of squares
is immediate. Now use $\sum y_j^2\le ND^2$ and
$\binom Nk\ge(N/k)^k$.
\end{proof}

For a polynomial $P$ of degree at most $N$, write $\operatorname{Pol}_R(P)$
for its symmetric multiaffine polarization on the $N$ variables indexed
by $R$. We have the exact expansion
\[
 \operatorname{Pol}_R(P)(x_j:j\in R)-P(x_R)
 =\sum_{k=2}^{\deg P}\frac{P^{(k)}(x_R)}{k!}
       \frac{e_k(x_j-x_R:j\in R)}{\binom Nk}.
\tag{3}
\]
The sum starts at two because the centered variables sum to zero.

\begin{lemma}[Integrated error]
For $d\ge16$ with $d\le c\sqrt m$ and $2d\le N$, where $c>0$ is a
sufficiently small absolute constant depending only on fixed $\theta$,
\begin{align*}
 (\ell+1)\E_w\E_B\bigl[A_B
   |\operatorname{Pol}_R(R_d)-R_d(x_R)|\bigr]&\le C d^4/m,\\
 (\ell+1)\E_w\E_B\bigl[A_B
   |\operatorname{Pol}_R(P_d)-P_d(x_R)|\bigr]&\le C d^2/m.
\end{align*}
\end{lemma}
\begin{proof}
The standard Laguerre bound
$|L_j^{(k)}(u)|\le\binom jk e^{u/2}$ for $u\ge0$ implies, by Leibniz
and Vandermonde's identity,
\[
 \frac{|(q_d(\theta u)^2)^{(k)}|}{k!}
 \le d^2 e^{\theta u}\frac{(4\theta d)^k}{(k!)^2}.
\tag{4}
\]
Indeed the Leibniz sum for $L_iL_j$ is at most
$2^k\sum_a\binom ia\binom j{k-a}=2^k\binom{i+j}k$.
For $R_d(u)=(2/(\theta d)-u)q_d(\theta u)^2$, differentiating its
linear factor in addition shows
\[
 \frac{|R_d^{(k)}(u)|}{k!}
 \le C_\theta d^2(1+u)e^{\theta u}
       \frac{(C_\theta d)^k(1+k^2)}{(k!)^2}.
\tag{5}
\]
For $P_d=h^{-2}(\sum_{j<h}L_j(\theta u))^2$, $h=\lfloor d/8\rfloor$,
the corresponding bound is
\[
 \frac{|P_d^{(k)}(u)|}{k!}
 \le C_\theta e^{\theta u}\frac{(C_\theta d)^k}{(k!)^2}.
\tag{6}
\]

At every row and for every $B$, (1) gives
$|x_j-x_R|\le C(s+1)$, while $x_R\le s/3$.
Insert (2), (5) into (3), multiply by $A_B$, and average over $B$.
Since $\E_B A_B\le e^{-\alpha s}$ and $\alpha\ge1/5$, the result
is bounded by
\[
 C d^2 e^{-c_0s}\sum_{k\ge2}
  \left(\frac{Cd}{\sqrt m}\right)^k
  \frac{k^{k/2}(1+k^2)}{(k!)^2}(1+s)^{k+1},
\tag{7}
\]
where $c_0>0$ is fixed. Extending the finite sum to infinity is harmless.

For each fixed nonnegative integer $a$, (1) also gives the explicit
factorial estimate
\[
 \E_w[(1+s)^{k+a}e^{-c_0s}]
 \le\frac{C^{k+a+1}(k+a)!}{m}.
\tag{8}
\]
To see this, bound $(1+s)^{k+a}e^{-c_0s/2}$ by
$C^{k+a}(k+a)!$ and apply the Laplace estimate to the remaining
exponential.
After multiplying (7) by $\ell+1=O(m)$ and using (8), the total error
is at most
\[
 C d^2\sum_{k\ge2}
   \left(\frac{Cd}{\sqrt m}\right)^k
   \frac{k^{k/2}(1+k^2)(k+1)!}{(k!)^2}.
\]
Since $k!\ge(k/e)^k$, the coefficient apart from the power is at
most $C^k(1+k)^3/\sqrt{k!}$. For $Cd/\sqrt m\le1$ the series
starting at $k=2$ is $O(d^2/m)$. This proves the first assertion.
The same argument with (6), and $a=0$, removes the prefactor $d^2$
and proves the second.
\end{proof}

The two factorials in (4) are essential to this estimate. Bounding
polynomial coefficients separately would instead introduce an
exponential factor in $d$ and recover only the earlier logarithmic
scale. This lemma supplies the polarization part of a prospective
power lower-bound proof; the exact beta identities, scalar test
positivity, and mixture comparison are separate inputs.

\section*{Variance-sensitive upgrade}
The stronger off-diagonal Gaussian bound proved in the main power
lower-bound note is
\[
 z^{h/2}|(L_j(\theta z))^{(h)}|
 \le(\theta j)^{h/2}e^{\theta z/2}\qquad(z>0).
\tag{9}
\]
Equivalently it follows from the modulus bound on the unitary Gaussian
translation matrix coefficient
$e^{-u/2}u^{h/2}\sqrt{(j-h)!/j!}\,L_{j-h}^{(h)}(u)$.
The following calculation is conditional only on that explicitly stated
classical inequality, in addition to (1).

\begin{lemma}[Variance-sensitive integrated polarization]
Let $Q_a(z)=\sum_{j=0}^d c_{aj}L_j(\theta z)$,
$B_a=\sum_j|c_{aj}|$, $a=1,2$, and put $P=Q_1Q_2$.
For $2d\le N$ and $d\le c m$, with a sufficiently small fixed $c>0$,
\[
 (\ell+1)\E_w\E_B\left[A_B
 |\operatorname{Pol}_R(P)-P(x_R)|\right]
 \le C B_1B_2\frac d m.
\tag{10}
\]
\end{lemma}
\begin{proof}
The proof of (2), before replacing the sum of squares by $ND^2$,
gives the stronger bound
\[
 \frac{|e_k(y)|}{\binom Nk}
 \le\left(\frac{\sqrt e\,\sigma\sqrt k}{N}\right)^k,
 \qquad \sigma^2=\sum_jy_j^2.
\tag{11}
\]
By (9) and the product rule, for $z>0$,
\[
 \frac{|P^{(k)}(z)|}{k!}
 \le B_1B_2e^{\theta z}
       \frac{(2\sqrt{\theta d/z})^k}{k!}.
\tag{12}
\]
At a fixed row and fixed $B$, write $v=x_R$ and $y_j=x_j-v$.
If $v=0$, all $x_j$ are zero, so the polarization error is zero.
Otherwise
\[
 \sigma^2\le\sum_{j\in R}x_j^2
 \le (\max_jx_j)\sum_{j\in R}x_j
 \le C N(s+1)v.
\tag{13}
\]
Insert (11)--(13) into the exact expansion (3). The powers of $v$
cancel. After averaging with $A_B$ and using (1), the result is bounded
by
\[
 C B_1B_2e^{-c_0s}\sum_{k=2}^{2d}
 \left(C\sqrt{d/m}\right)^k
      \frac{k^{k/2}}{k!}(s+1)^{k/2}.
\]
The factorial weighted-moment estimate (8), with
$\lceil k/2\rceil$ in place of $k+a$, shows that the normalized
integrated error is at most
\[
 C B_1B_2\sum_{k=2}^{2d}
 \left(C\sqrt{d/m}\right)^k
       \frac{k^{k/2}\lceil k/2\rceil!}{k!}.
\]
For even and odd $k$ separately, the elementary factorial bounds
$(k/e)^k\le k!\le k^k$ give
\[
 \frac{k^{k/2}\lceil k/2\rceil!}{k!}
 \le C^k(k+1).
\]
Thus the series is geometric with a polynomial factor, starts at
$k=2$, and is $O(B_1B_2d/m)$ whenever $d\le c m$ with $c$ sufficiently
small. This proves (10).
\end{proof}

For the endpoint taper, the recurrence identity
$yq_d(y)=d^{-1}\sum_{j<d}L_j(y)-L_d(y)$ factors
$R_d=q_d(\theta z)[(2/(\theta d)-z)q_d(\theta z)]$ into two Laguerre
expansions with coefficient sums $O(d)$ and $O(1)$.
Hence (10) bounds its polarization error by $O(d^2/m)$.
This is compatible with endpoint degree $d$ of order $\sqrt m$;
the weighted-mixture comparison must satisfy the corresponding bound
as a separate input.

\end{document}
